\chapter{Facility and instrument}
\label{chap:setup}

\section{The venue}

TV2 is a television studio of the PUCP Faculty of Communications, borrowed for
the session through internal services. The project does not own a space that
can be made dark at this scale, and the studio offers three things the
experiment needs and the laboratory enclosure does not: a room that can be
blacked out completely, professional lighting fixtures with the rigging to
position and aim them, and enough floor depth to place a fixture far from the
subject, which is what the collimation requirement A2.3 ultimately demands.

A studio also brings two hazards that the acceptance criteria are written to
catch. Studio lighting is designed for an aesthetic result, so fill light,
bounce and multiple sources are the norm and have to be deliberately
suppressed, which is criterion A2.2. And studio fixtures are frequently
dimmed, often by a method whose modulation frequency is unknown, which is
criterion A2.7.

\section{The instrument, as built}

\begin{table}[htbp]
\centering
\caption[Instrument configuration]{The camera payload as assembled for this
experiment. The integrated support was completed and verified working on 13
October 2026; it is the geometry the calibration must be performed on.}
\label{tab:instrument}
\begin{tabularx}{\textwidth}{@{}l X@{}}
\toprule
\textbf{Item} & \textbf{Configuration} \\
\midrule
Compute & Raspberry Pi Compute Module 5, Cortex-A76 at 2.4~GHz, 16~GB \\
Sensors & Two Arducam HQ 477 modules, Sony IMX477, 12.3~MP
          \cite{sony_imx477_datasheet} \\
Lenses & Arducam C1508ZM04, 8~mm, f/1.6 to f/16, C-mount
         \cite{arducam_c1508zm04} \\
Baseline & 100~mm, fixed by the integrated support \\
Capture & $4056 \times 3040$, 12-bit sensor mode, explicit
          \texttt{-{}-mode} argument \\
Preview & $1920 \times 1440$, used for guidance only, never for measurement \\
Frame rate & 15~frame/s in the current software configuration \\
Synchronisation & Sequential software triggering, with the inter-camera delay
                  recorded per pair \\
Calibration target & ChArUco $5 \times 7$, 37~mm squares, 29.6~mm markers,
                     \texttt{DICT\_6X6\_250} \cite{garrido2014aruco} \\
\bottomrule
\end{tabularx}
\end{table}

\textbf{On the calibration target.} It is a ChArUco pattern printed on paper
and mounted on a rigid folder. Two properties of a target matter and neither
is usually verified: flatness and detectability. Detectability is the reason
for ChArUco rather than a plain chessboard, since a self-identifying target
labels every corner individually and tolerates partial views, which is what
makes interactive guidance possible \cite{rojtberg2018pose}. Flatness has
never been measured for this board, and a paper sheet on a folder is the
weaker of the two properties. It is recorded here as a known limitation
rather than as a controlled parameter.

\textbf{On synchronisation.} Hardware triggering is not exposed through the
standard development configuration, so the two cameras are triggered
sequentially in software. The camera system report asserts this is adequate
for target motion below roughly 5~cm/s \cite{linku_d5_camera_report}; that
bound is an assumption rather than a measured synchronisation error. Hand
motion will exceed it. The inter-camera delay is recorded per pair, and pairs
whose delay is large are candidates for exclusion from A3.5. For reference,
the February 2026 reference pair records a 10.4~ms delay.

\section{The subject}

A physical tether proxy, nominally 1~mm in diameter, strung across the LINKU
structure. The structure itself is built by another member of the team and is
outside the scope of this document; what matters here is that the tether is
attached to it at two points whose positions are recorded, and that the
camera support is mounted on it.

The 1~mm figure is the physical proxy actually used in the laboratory. It
should not be confused with the 0.8~mm figure that appears in the camera
system report, which is explicitly described there as a conservative
assumption adopted for the analytical detectability estimate and is not a
measurement of anything \cite{linku_d5_camera_report}.

\section{Working distance and what the optics do at it}

The nominal working distance is 3.5~m, taken from the Blender replica of the
structure. At that distance, with the fitted 8~mm lens on the IMX477, the
tool \cite{linku_dof_visualizer} gives the figures in Table
\ref{tab:optics-at-work}.

\begin{table}[htbp]
\centering
\caption[Optical behaviour at the working distance]{Optical behaviour of the
fitted configuration at the nominal 3.5~m working distance, computed for the
IMX477 with a 1.55~\textmu m pixel and a one-pixel circle of confusion. These
figures do not transfer to the IMX708 B0308 flight sensor, which has a
different pixel pitch and a different focal length.}
\label{tab:optics-at-work}
\begin{tabularx}{\textwidth}{@{}X l l@{}}
\toprule
\textbf{Quantity} & \textbf{Value} & \textbf{Note} \\
\midrule
Projected width of a 1~mm tether at 3.5~m & 1.47~px & marginally resolved \\
Depth of field focused at 3.5~m, f/1.6 & 3.08 to 4.05~m & 0.96~m total \\
Hyperfocal distance at f/1.6 & 25.8~m & \\
Depth of field focused at 21~m & 11.58 to 112.6~m & tether at 3.5~m is blurred \\
Airy disc diameter at f/1.6 & 2.15~\textmu m & 1.39~px \\
Maximum exposure at the 1\textdegree/s reference rate & 11.1~ms & 1~px blur budget \\
Minimum illuminance to saturate a white card in that exposure & 2320~lx & criterion A2.5 \\
\bottomrule
\end{tabularx}
\end{table}

Two consequences drive the procedure in Chapter \ref{chap:procedure}.

\textbf{The depth of field is under a metre.} Focusing for a far subject does
not merely reduce sharpness at 3.5~m, it moves the near limit of the sharp
zone out to 11.6~m and leaves the tether entirely outside it. The focus has to
be set by looking at a live image at full resolution, at the real distance,
and then locked.

\textbf{The focus ring carries no usable scale.} Its full travel moves the
lens by 0.45~mm, and the last 0.003~mm of that, 0.7\% of the rotation, covers
everything from 21~m to infinity. Setting focus by reading a number off the
barrel is not a viable method on this lens.
